% =============================================================================
% Apêndice — Resultados experimentais complementares
% =============================================================================
\chapter{Resultados experimentais complementares}
\label{ap:experimentos}

O Capítulo~5 apresenta em detalhe três experimentos --- a linha de base com o conjunto completo de \textit{features} (Experimento~1), a configuração enxuta de 11 \textit{features} usada como referência operacional (Experimento~2) e o teste de generalização exclusivamente em curvas reais (Experimento~3) ---, além de sintetizar todos os resultados na Tabela~\ref{tab:comparacao_f1}. Este apêndice reúne os \textbf{experimentos intermediários} cujas conclusões já foram incorporadas àquela síntese, mas que se documentam aqui por completude: a redução progressiva do conjunto de \textit{features} para 14 (Experimento~4) e 13 variáveis (Experimento~5), a ablação isolada da \textit{feature} \texttt{kmeans\_centroid\_dist} (Experimento~6) e uma variante do Experimento~3 com divisão 65\%/35\%, utilizada como ablação do tamanho do conjunto de treino. Todas as métricas referem-se ao conjunto de teste e seguem o protocolo descrito na Seção~\ref{sec:config_experimental}.

% -----------------------------------------------------------------------------
\section{Experimento 4: subconjunto de 14 \textit{features}}
\label{sec:resultados_exp2}

No Experimento 4 foi utilizado um subconjunto reduzido de \textbf{14 \textit{features}}, removendo amplitude, desvio padrão do fluxo bruto, máximo do Savitzky-Golay, todas as derivadas e duas métricas observacionais (\texttt{Occ\_flux\_min} e \texttt{Occ\_flux\_min\_over\_baseline}). O objetivo é diminuir o risco de \textit{overfitting} e melhorar a razão entre número de variáveis e tamanho do conjunto de treino. A Tabela~\ref{tab:metricas_2} apresenta os resultados.

\begin{table}[htbp]
\centering
\caption{Métricas de desempenho no conjunto de teste --- Experimento 4. Subconjunto de 14 \textit{features}, divisão 80\%/20\%.}
\label{tab:metricas_2}
\begin{tabular}{lccccc}
\hline
\textbf{Modelo} & \textbf{Acurácia} & \textbf{Precisão} & \textbf{Sensibilidade} & \textbf{F1-score} & \textbf{AUC-ROC} \\
\hline
CatBoost        & \textbf{0,9912} & 0,9937 & 0,9875 & \textbf{0,9906} & 0,9996 \\
Random Forest   & 0,9882 & 0,9937 & 0,9813 & 0,9874 & 0,9997 \\
XGBoost         & 0,9882 & 0,9875 & 0,9875 & 0,9875 & 0,9996 \\
Regressão Log.  & 0,9882 & 1,0000 & 0,9750 & 0,9873 & 1,0000 \\
\hline
\end{tabular}
\end{table}

O melhor F1-score foi alcançado pelo \textbf{CatBoost} (F1 = 0,9906). As métricas permanecem muito altas, indicando que a redução de 28 para 14 \textit{features} preserva a informação discriminativa. A Figura~\ref{fig:confusao_2} apresenta as matrizes de confusão.

\begin{figure}[htbp]
\begin{center}
	\includegraphics[width=0.85\textwidth]{Tese/pngs/exp2/confusion_matrices.png}
\caption{Matrizes de confusão dos quatro modelos no Experimento 4 (14 \textit{features}).}
\label{fig:confusao_2}
\end{center}
\end{figure}

% -----------------------------------------------------------------------------
\section{Experimento 5: 13 \textit{features} (sem \texttt{Feature\_Savgol\_Min})}
\label{sec:resultados_exp3}

No Experimento 5 foi removida a \textit{feature} \texttt{Feature\_Savgol\_Min}, mantendo 13 \textit{features}. O objetivo é avaliar a contribuição dessa variável, que costuma figurar entre as mais relevantes nos gráficos de importância. A Tabela~\ref{tab:metricas_3} resume as métricas.

\begin{table}[htbp]
\centering
\caption{Métricas de desempenho no conjunto de teste --- Experimento 5. 13 \textit{features} (sem \texttt{Feature\_Savgol\_Min}), divisão 80\%/20\%.}
\label{tab:metricas_3}
\begin{tabular}{lccccc}
\hline
\textbf{Modelo} & \textbf{Acurácia} & \textbf{Precisão} & \textbf{Sensibilidade} & \textbf{F1-score} & \textbf{AUC-ROC} \\
\hline
Random Forest   & \textbf{0,9912} & 1,0000 & 0,9813 & \textbf{0,9905} & 0,9998 \\
CatBoost        & 0,9882 & 1,0000 & 0,9750 & 0,9873 & 0,9998 \\
Regressão Log.  & 0,9882 & 1,0000 & 0,9750 & 0,9873 & 1,0000 \\
XGBoost         & 0,9853 & 0,9874 & 0,9813 & 0,9843 & 0,9997 \\
\hline
\end{tabular}
\end{table}

O \textbf{Random Forest} lidera com F1 = 0,9905. Em relação ao Experimento 4, a remoção de \texttt{Feature\_Savgol\_Min} causa queda moderada em alguns modelos, mas o desempenho global permanece elevado, indicando que as demais \textit{features} absorvem parte da informação. As matrizes de confusão estão na Figura~\ref{fig:confusao_3}. Vale notar que a análise de ajuste de limiar e a validação cruzada apresentadas no Capítulo~5 (Seções~\ref{sec:threshold_tuning} e~\ref{sec:cross_validation}) foram conduzidas sobre esta configuração de 13 \textit{features}.

\begin{figure}[htbp]
\begin{center}
	\includegraphics[width=0.85\textwidth]{Tese/pngs/exp3/confusion_matrices.png}
\caption{Matrizes de confusão dos quatro modelos no Experimento 5 (13 \textit{features}, sem \texttt{Feature\_Savgol\_Min}).}
\label{fig:confusao_3}
\end{center}
\end{figure}

% -----------------------------------------------------------------------------
\section{Experimento 6: ablação isolada do K-Means}
\label{sec:resultados_exp4}

No Experimento 6 foi removida \texttt{kmeans\_centroid\_dist}, mantendo 12 \textit{features} (sem \texttt{Feature\_Savgol\_Min}), com divisão 80\%/20\%. A pergunta investigada é se a alta importância atribuída a essa variável pelos modelos de árvore reflete uma contribuição insubstituível ou se as demais \textit{features} absorvem sua capacidade preditiva. A Tabela~\ref{tab:metricas_4} apresenta as métricas.

\begin{table}[htbp]
\centering
\caption{Métricas de desempenho no conjunto de teste --- Experimento 6. 12 \textit{features} (sem \texttt{kmeans\_centroid\_dist}), divisão 80\%/20\%.}
\label{tab:metricas_4}
\begin{tabular}{lccccc}
\hline
\textbf{Modelo} & \textbf{Acurácia} & \textbf{Precisão} & \textbf{Sensibilidade} & \textbf{F1-score} & \textbf{AUC-ROC} \\
\hline
Random Forest   & \textbf{0,9912} & 0,9937 & 0,9875 & \textbf{0,9906} & 0,9998 \\
XGBoost         & \textbf{0,9912} & 0,9937 & 0,9875 & \textbf{0,9906} & 0,9997 \\
CatBoost        & 0,9882 & 0,9937 & 0,9813 & 0,9874 & 0,9997 \\
Regressão Log.  & 0,9882 & 1,0000 & 0,9750 & 0,9873 & 1,0000 \\
\hline
\end{tabular}
\end{table}

Random Forest e XGBoost empatam com F1 = 0,9906 --- desempenho igual ou superior ao da configuração que mantém o K-Means. A Figura~\ref{fig:confusao_4} apresenta as matrizes de confusão.

\begin{figure}[htbp]
\begin{center}
	\includegraphics[width=0.85\textwidth]{Tese/pngs/exp4/confusion_matrices.png}
\caption{Matrizes de confusão dos quatro modelos no Experimento 6 (12 \textit{features}, sem \texttt{kmeans\_centroid\_dist}).}
\label{fig:confusao_4}
\end{center}
\end{figure}

A Tabela~\ref{tab:ablation_comparacao} confronta os F1-scores do Experimento~5 (13 \textit{features}, com \texttt{kmeans}), do Experimento~6 (12 \textit{features}, sem \texttt{kmeans}) e do Experimento~2 (11 \textit{features}, sem \texttt{Feature\_Savgol\_Min} e sem \texttt{kmeans}), evidenciando que a remoção da \textit{feature} de K-Means não degrada o desempenho global.

\begin{table}[htbp]
\centering
\caption{Comparação do F1-score: com \texttt{kmeans} (Exp.~5), sem \texttt{kmeans} (Exp.~6) e ablação combinada (Exp.~2). Divisão 80\%/20\%.}
\label{tab:ablation_comparacao}
\begin{tabular}{lccc}
\hline
\textbf{Modelo} & \textbf{Exp.~5 (13 feat.)} & \textbf{Exp.~6 (12 feat., s/ kmeans)} & \textbf{Exp.~2 (11 feat.)} \\
\hline
Random Forest   & \textbf{0,9905} & \textbf{0,9906} & 0,9874 \\
XGBoost         & 0,9843 & \textbf{0,9906} & \textbf{0,9906} \\
CatBoost        & 0,9873 & 0,9874 & \textbf{0,9905} \\
Regressão Log.  & 0,9873 & 0,9873 & 0,9841 \\
\hline
\end{tabular}
\end{table}

A importância antes concentrada em \texttt{kmeans\_centroid\_dist} é redistribuída entre variáveis como \texttt{Feature\_Savgol\_std}, \texttt{Occ\_depth}, \texttt{Occ\_duration\_s} e \texttt{MinLogP\_Ttest}. Esse resultado sustenta a conclusão discutida no Capítulo~5: \texttt{kmeans\_centroid\_dist} e \texttt{Feature\_Savgol\_std} capturam o mesmo fenômeno físico (a dispersão do fluxo causada pelos dois patamares dentro e fora do \textit{dip}), de modo que a alta importância atribuída ao K-Means não reflete informação insubstituível. Na prática, isso permite simplificar a \textit{pipeline} eliminando o passo de K-Means sem perda de desempenho.

Por fim, a Figura~\ref{fig:confusao_5} apresenta as matrizes de confusão da configuração enxuta de 11 \textit{features} (Experimento~2), discutida em detalhe na Seção~\ref{sec:resultados_exp5} do Capítulo~5.

\begin{figure}[htbp]
\begin{center}
	\includegraphics[width=0.85\textwidth]{Tese/pngs/exp5/confusion_matrices.png}
\caption{Matrizes de confusão dos quatro modelos no Experimento 2 (11 \textit{features}).}
\label{fig:confusao_5}
\end{center}
\end{figure}

% -----------------------------------------------------------------------------
\section{Variante do Experimento 3 com divisão 65/35: ablação do tamanho do treino}
\label{sec:exp3_6535}

O Experimento~3 do Capítulo~5 avalia a generalização para curvas exclusivamente reais sob a mesma divisão 80\%/20\% dos demais experimentos. Esta seção documenta uma variante com divisão 65\%/35\%, que cumpre dois papéis: (i)~ampliar o conjunto de teste real avaliado (347 curvas: 281 positivas e 66 negativas, contra 198 na configuração principal), reduzindo a incerteza amostral das métricas; e (ii)~servir de \textbf{ablação do tamanho do conjunto de treino}. Ao remover cerca de 15\% dos dados de treinamento, mede-se a sensibilidade do desempenho à quantidade de exemplos disponíveis.

\begin{table}[htbp]
\centering
\caption{Métricas de desempenho no conjunto de teste --- variante do Experimento 3 com divisão 65\%/35\% (teste somente em curvas reais: 347 curvas, 281 positivas e 66 negativas).}
\label{tab:metricas_exp3_6535}
\begin{tabular}{lccccc}
\hline
\textbf{Modelo} & \textbf{Acurácia} & \textbf{Precisão} & \textbf{Sensibilidade} & \textbf{F1-score} & \textbf{AUC-ROC} \\
\hline
XGBoost         & \textbf{0,9741} & 0,9928 & 0,9751 & \textbf{0,9838} & 0,9968 \\
Random Forest   & 0,9683 & 0,9856 & 0,9751 & 0,9803 & 0,9956 \\
CatBoost        & 0,9654 & 0,9891 & 0,9680 & 0,9784 & 0,9967 \\
Regressão Log.  & 0,9597 & 0,9926 & 0,9573 & 0,9746 & 0,9959 \\
\hline
\end{tabular}
\end{table}

A comparação com a configuração principal (Tabela~\ref{tab:metricas_7}, Capítulo~5) é direta: com 15\% menos dados de treino, o melhor F1 recua de 0,9906 para 0,9838 (queda de apenas $\approx$0,7~ponto percentual) e a AUC-ROC permanece $\geq 0{,}9956$ em todos os modelos. Essa estabilidade indica que a curva de aprendizado dos modelos está \textbf{saturada} na região de operação: o desempenho depende fracamente do volume de treino, comportamento incompatível com sobreajuste dominante (um modelo sobreajustado degradaria de forma acentuada ao perder dados de ajuste). O resultado converge com a validação cruzada estratificada (Seção~\ref{sec:cross_validation}), que apresenta desvios padrão pequenos entre dobras. As Figuras~\ref{fig:confusao_7} e~\ref{fig:ROC7} apresentam as matrizes de confusão e as curvas ROC desta variante.

\begin{figure}[htbp]
\begin{center}
	\includegraphics[width=0.85\textwidth]{Tese/pngs/exp7/confusion_matrices.png}
\caption{Matrizes de confusão dos quatro modelos na variante 65/35 do Experimento 3 (teste somente em curvas reais).}
\label{fig:confusao_7}
\end{center}
\end{figure}

\begin{figure}[htbp]
\begin{center}
	\includegraphics[width=0.7\textwidth]{Tese/pngs/exp7/roc_curves.png}
\caption{Curvas ROC --- variante 65/35 do Experimento 3 (teste somente em curvas reais).}
\label{fig:ROC7}
\end{center}
\end{figure}

% -----------------------------------------------------------------------------
\section{Curvas ROC dos experimentos complementares}
\label{sec:roc_complementares}

As curvas ROC do Experimento 1 são apresentadas no Capítulo~5 (Figura~\ref{fig:ROC1}); as da variante 65/35 do Experimento~3 estão na Seção~\ref{sec:exp3_6535} (Figura~\ref{fig:ROC7}). Para completude, as Figuras~\ref{fig:ROC2}--\ref{fig:ROC5} reúnem as curvas ROC dos demais experimentos (2, 4, 5 e 6). Em todos os casos a AUC permanece acima de 0,999, confirmando a elevada separabilidade das classes no conjunto de teste misto.

\begin{figure}[htbp]
\begin{center}
	\includegraphics[width=0.7\textwidth]{Tese/pngs/exp2/roc_curves.png}
\caption{Curvas ROC --- Experimento 4 (14 \textit{features}).}
\label{fig:ROC2}
\end{center}
\end{figure}

\begin{figure}[htbp]
\begin{center}
	\includegraphics[width=0.7\textwidth]{Tese/pngs/exp3/roc_curves.png}
\caption{Curvas ROC --- Experimento 5 (13 \textit{features}).}
\label{fig:ROC3}
\end{center}
\end{figure}

\begin{figure}[htbp]
\begin{center}
	\includegraphics[width=0.7\textwidth]{Tese/pngs/exp4/roc_curves.png}
\caption{Curvas ROC --- Experimento 6 (12 \textit{features}, sem \texttt{kmeans}).}
\label{fig:ROC4}
\end{center}
\end{figure}

\begin{figure}[htbp]
\begin{center}
	\includegraphics[width=0.7\textwidth]{Tese/pngs/exp5/roc_curves.png}
\caption{Curvas ROC --- Experimento 2 (11 \textit{features}).}
\label{fig:ROC5}
\end{center}
\end{figure}

% -----------------------------------------------------------------------------
\section{Importância das \textit{features} nos Experimentos 4 e 5}
\label{sec:importancia_complementar}

O Capítulo~5 discute a importância das variáveis com base na figura do Experimento~1 (conjunto completo). As Figuras~\ref{fig:feature_importance_2} e~\ref{fig:feature_importance_3} apresentam os gráficos análogos para os Experimentos 4 e 5, úteis para acompanhar como a relevância relativa das variáveis se redistribui à medida que o conjunto é reduzido. A interpretação dos valores deve considerar as diferenças metodológicas entre os métodos de cada modelo, discutidas na Seção~\ref{sec:feature_importance_metodologia}.

\begin{figure}[p]
\centering
\noindent\textbf{Experimento 4 (14 \textit{features}).}\\[0.5em]
\begin{minipage}[t]{0.48\textwidth}
	\centering
	\includegraphics[width=\linewidth]{Tese/pngs/exp2/feature_importance_logistic_regression.png}
	\vspace{0.3em}
	\small (a) Regressão Log.
\end{minipage}\hfill
\begin{minipage}[t]{0.48\textwidth}
	\centering
	\includegraphics[width=\linewidth]{Tese/pngs/exp2/feature_importance_random_forest.png}
	\vspace{0.3em}
	\small (b) Random Forest
\end{minipage}\\[1em]
\begin{minipage}[t]{0.48\textwidth}
	\centering
	\includegraphics[width=\linewidth]{Tese/pngs/exp2/feature_importance_xgboost.png}
	\vspace{0.3em}
	\small (c) XGBoost
\end{minipage}\hfill
\begin{minipage}[t]{0.48\textwidth}
	\centering
	\includegraphics[width=\linewidth]{Tese/pngs/exp2/feature_importance_catboost.png}
	\vspace{0.3em}
	\small (d) CatBoost
\end{minipage}\\[0.5em]
\caption{Importância de variáveis por modelo --- Experimento 4 (14 \textit{features}). (a) Regressão Logística, (b) Random Forest, (c) XGBoost, (d) CatBoost.}
\label{fig:feature_importance_2}
\end{figure}

\begin{figure}[p]
\centering
\noindent\textbf{Experimento 5 (13 \textit{features}, sem \texttt{Feature\_Savgol\_Min}).}\\[0.5em]
\begin{minipage}[t]{0.48\textwidth}
	\centering
	\includegraphics[width=\linewidth]{Tese/pngs/exp3/feature_importance_logistic_regression.png}
	\vspace{0.3em}
	\small (a) Regressão Log.
\end{minipage}\hfill
\begin{minipage}[t]{0.48\textwidth}
	\centering
	\includegraphics[width=\linewidth]{Tese/pngs/exp3/feature_importance_random_forest.png}
	\vspace{0.3em}
	\small (b) Random Forest
\end{minipage}\\[1em]
\begin{minipage}[t]{0.48\textwidth}
	\centering
	\includegraphics[width=\linewidth]{Tese/pngs/exp3/feature_importance_xgboost.png}
	\vspace{0.3em}
	\small (c) XGBoost
\end{minipage}\hfill
\begin{minipage}[t]{0.48\textwidth}
	\centering
	\includegraphics[width=\linewidth]{Tese/pngs/exp3/feature_importance_catboost.png}
	\vspace{0.3em}
	\small (d) CatBoost
\end{minipage}\\[0.5em]
\caption{Importância de variáveis por modelo --- Experimento 5 (13 \textit{features}, sem \texttt{Feature\_Savgol\_Min}). (a) Regressão Logística, (b) Random Forest, (c) XGBoost, (d) CatBoost.}
\label{fig:feature_importance_3}
\end{figure}
